\documentclass{situationist-preprint}

\title{Finitude and Meaning}
\msid{EXI-19-b}
\msfamily{EXI-19}
\msdiscipline{Meaning, Death and the Absurd}
\mstype{Companion manuscript}
\mskeywords{finitude; meaning in life; mortality; shape of a life; scarcity; immortality}

\begin{document}
\maketitle

\begin{abstract}
It is often said that finitude is the source of meaning: that our projects matter because our time is limited, and that an immortal being could find nothing worth doing. EXI-19-a argues that death has no meaning in itself but gives life a shape by limiting it. We draw out what this implies about finitude and meaning. Finitude is not a source of meaning, since it adds no object of worth to engage with. It is a condition of one kind of meaning, that of a life as a whole, which depends on the life's shape. And, combined with the result of EXI-03, it is a condition of a life's remaining meaningful throughout. The scarcity thesis of Nussbaum and Scheffler is right about the second and third points and overstated as a claim about the first.
\end{abstract}

\section{Three claims about finitude}

On the account of EXI-01 and EXI-17, a life is meaningful to the extent that its subject actively engages with what is worthwhile. Finitude might bear on meaning in three ways: as a source of it, as a condition of the meaning of a life as a whole, or as a condition of the persistence of engagement. We take them in turn.

\begin{proposition}[Not a source]\label{p:source}
That a life is limited by death is not something whose engagement makes the life meaningful, and does not make anything engaged with more worthwhile.
\end{proposition}

\begin{proof}
By Proposition~3.1 of EXI-19-a, death has no meaning in itself. The worth of the objects of engagement, knowledge, persons, beauty and the good of others, does not depend on how long those who engage them live: a proof is no less a proof for being understood by someone with unlimited time. So the limit adds neither an object of worth nor worth to an object, and by the account of EXI-01 it is not a source of meaning.
\end{proof}

\begin{proposition}[Whole-life meaning]\label{p:whole}
In so far as the meaning of a life as a whole depends on its shape, finitude is a condition of it.
\end{proposition}

\begin{proof}
Some meaning belongs to a life only as a whole: that it was a life of growing understanding, of a commitment kept, of a wrong redeemed. Such meaning depends on the temporal order of the engagements it contains, that is, on the life's shape (Definition~2.2 of EXI-19-a; \citealp{Velleman1991}). By Assumption~2.3 of EXI-19-a, a sequence has an order as a whole only if it has a last member, and by Theorem~3.2 of EXI-19-a it is death that supplies it. So the meaning that a life has as a whole requires that it be limited.
\end{proof}

\begin{proposition}[Persistence of engagement]\label{p:persist}
Finitude is a condition of a life's remaining meaningful throughout its course.
\end{proposition}

\begin{proof}
By the result of EXI-03, an endless life would eventually exhaust the categorical desires that give a life its point \citep{Williams1973}. A life without categorical desires lacks the active engagement that meaning requires (EXI-17). So an endless life would at some point cease to be meaningful, and a life remains meaningful throughout its course only if it ends before this point.
\end{proof}

Proposition~\ref{p:persist} concerns the eventual exhaustion of engagement, not its quality at any given time. An immortal's engagement with inquiry in the first thousand years would be engagement with what is worthwhile, as meaningful as ours; the claim is that it would not go on being so.

\section{The scarcity thesis}

\citet{Nussbaum1989} argues that many human values and virtues presuppose mortality, and \citet{Scheffler2013} that our capacity to value depends on the awareness of temporal scarcity, so that an immortal life would lack the structure in which valuing makes sense.

\begin{objection}[Scarcity]
Projects matter to us because time is limited: choosing one is forgoing others. Without scarcity nothing would be urgent, and nothing would matter. So finitude is a source of meaning after all.
\end{objection}

\begin{reply}
Scarcity gives choices their cost and projects their urgency, and urgency affects how we engage. But Proposition~\ref{p:source} concerns the worth of what we engage with, and urgency adds no worth: the friend's good matters no less to someone who will have other friends. Where the scarcity thesis is right, it is right about Propositions~\ref{p:whole} and \ref{p:persist}. Scheffler's point that an immortal life would lack the stages of a human life is a point about shape, and Nussbaum's point that some values belong to a life of a certain form is a point about the conditions under which those engagements persist. Neither shows that the limit is something whose engagement confers meaning.
\end{reply}

\begin{objection}[A benefit after all]
If finitude is a condition of whole-life meaning and of the persistence of engagement, then death benefits the one who dies, contrary to EXI-02.
\end{objection}

\begin{reply}
It does not follow. The conditions are met by a limit, not by any particular one; a later death would also have given the life a shape and would, in the ordinary case, have left more time for engagement, as Corollary~3.3 of EXI-19-a notes \citep{Fischer1994}. A death that deprives its subject of goods remains bad for her, though some limit is a condition of meaning.
\end{reply}

\section{Conclusion}

Finitude does not make life meaningful. It makes possible the meaning of a life as a whole, through the shape that a limit gives, and it keeps the life within the span in which engagement persists. That is the sense in which death, having no meaning in itself, bears on the meaning of the life it ends.

\begin{thebibliography}{9}\small
\bibitem[Fischer(1994)]{Fischer1994} Fischer, J.~M. (1994). Why immortality is not so bad. \emph{International Journal of Philosophical Studies}, 2(2), 257--270.
\bibitem[Nussbaum(1989)]{Nussbaum1989} Nussbaum, M.~C. (1989). Mortal immortals: Lucretius on death and the voice of nature. \emph{Philosophy and Phenomenological Research}, 50(2), 303--351.
\bibitem[Scheffler(2013)]{Scheffler2013} Scheffler, S. (2013). \emph{Death and the Afterlife}. Oxford University Press.
\bibitem[Velleman(1991)]{Velleman1991} Velleman, J.~D. (1991). Well-being and time. \emph{Pacific Philosophical Quarterly}, 72(1), 48--77.
\bibitem[Williams(1973)]{Williams1973} Williams, B. (1973). The Makropulos case: Reflections on the tedium of immortality. In \emph{Problems of the Self}, 82--100. Cambridge University Press.
\end{thebibliography}

\end{document}
